We run two families of experiments: a prompt-component ablation on single-kernel benchmarks, and an
end-to-end study on production applications.

\paragraph{Benchmarks.}
The prompt ablation uses six HeCBench \cite{jin2023hecbench} CUDA benchmarks---\texttt{atan2},
\texttt{bitonic-sort}, \texttt{md5hash}, \texttt{accuracy}, \texttt{sobol} and
\texttt{shmembench}---chosen to span memory-bound and compute-bound behaviour and to include kernels
with built-in verification. The profile-data ablation (Section~\ref{sec:profabl}) uses a further eleven
HeCBench applications: \texttt{adv}, \texttt{amgmk}, \texttt{atomicSystemWide},
\texttt{background-subtract}, \texttt{bincount}, \texttt{blockAccess}, \texttt{boxfilter},
\texttt{ccsd-trpdrv}, \texttt{chemv}, \texttt{concat} and \texttt{convolutionSeparable}.

\paragraph{Applications.}
The end-to-end study uses three production codes: AutoDock-GPU (molecular docking, CUDA),
LAMMPS (molecular dynamics, Kokkos/CUDA) and Cholla (astrophysical hydrodynamics, CUDA). Each is built
from its own build system with its own inputs, and validated by a gate derived from its scientific
output rather than from a benchmark harness (Section~\ref{sec:realapps}).

\paragraph{Models.}
The reported prompt-ablation grid uses GPT-5 and Gemini~2.5~Pro. The application and profile-data
studies use GPT-5 and Gemini~3.1~Pro; we note the model change explicitly wherever numbers from the two
studies appear together, since they are not directly comparable.
% Earlier exploratory runs used llama3.1:70B, llama2:7B and qwen-coder:14B. They are not in any
% reported table; add a paragraph here if you want to include those numbers.

\paragraph{Hardware.}
All benchmark runs---the prompt ablation and the HeCBench part of the profile-data ablation---are
generated and timed on an NVIDIA H100 PCIe (sm\_90). The application study and the real-application
part of the profile-data ablation run on an NVIDIA GH200 (sm\_90). Both nodes use CUDA 12.6.3. One
mismatch is worth stating: the profiles supplied in the prompt for the six original ablation benchmarks
were collected on an A100 ($108$ SMs) rather than on the H100 ($114$ SMs) they were then timed on; the
eleven profile-data-ablation benchmarks use native H100 profiles. All nodes are shared with other
users, which motivates the measurement protocol of Section~\ref{sec:method}: build once, interleave
variants within one session, wait for an idle GPU, and log competing GPU processes at the start and end
of every repetition so that contaminated repetitions can be excluded rather than averaged in.
